
\exo

\begin{enumerate}
\item La fonction $f$ est définie sur $\R$ par $f(x)=3x^2+7x$.

Déterminer les images par $f$ des nombres suivants :
\[2\qpvq -3\qpvq 0\qpvq \sqrt{5}.\]
\item La fonction $g$ est définie sur $\R$ par $g(x)=\frac{4}{3}x+5$.

\begin{enumerate}
\item Calculer $g(6)$ et $g(7)$.
\item Quelle est l'image de $-5$ par $g$ ?
\end{enumerate}
\end{enumerate}

\exo

\begin{enumerate}
\item La fonction $f$ est définie sur $\R$ par $f(x)=3x-8$.

Déterminer les éventuels antécédents par $f$ des nombres suivants :
\[3\qpvq -5\qpvq \tfrac{1}{2}\qpvq 0,1.\]
\item La fonction $g$ est définie sur $\R$ par $g(x)=(x-5)(3-2x)$.

Déterminer les éventuels antécédents de $0$ par $g$.
\end{enumerate}

\exo

La fonction $f$ est définie sur $\R$ par $f(x)=\dfrac{4x+2}{x^2+1}$.
\begin{enumerate}
\item Est-il vrai que $3$ est un antécédent de $1$ par $f$ ?
\item Est-il vrai que $0$ et $2$ ont la même image par $f$ ?
\item Déterminer l'image de $\frac{1}{2}$ par $f$.
\item Déterminer les éventuels antécédents de $0$ par $f$.
\end{enumerate}

\exo

\vspace{-12pt}
\setlength{\columnsep}{1cm}
\setlength{\columnseprule}{0pt}
\begin{multicols}{2}
La courbe ci-contre est la représentation graphique d'une fonction $f$ définie sur $\R$.

Par lectures graphiques, déterminer :
\begin{enumerate}
\item l'image de $-1$ par $f$ ;
\item l'image de $0$ par $f$ ;
\item les éventuels antécédents de $1$ par $f$ ;
\item les éventuels antécédents de $3$ par $f$.
\end{enumerate}

\begin{center}
\def\xmin{-4}
\def\xmax{4}
\def\ymin{-0.5}
\def\ymax{4}
\def\xunite{0.9cm}
\def\yunite{0.9cm}
\psset{unit=\xunite, xunit=\xunite, yunit=\yunite, algebraic=true, dimen=middle, dotstyle=*, dotsize=3pt 0, linewidth=0.8pt, arrowsize=3pt 2, arrowinset=0.25, comma=true}
\begin{pspicture*}(\xmin,\ymin)(\xmax,\ymax)
\psgrid[xunit=\xunite, yunit=\yunite, subgriddiv=0, gridlabels=0, gridcolor=lightgray](0,0)(\xmin,\ymin)(\xmax,\ymax)
\psaxes[labelFontSize=\scriptstyle, xAxis=true, yAxis=true, Dx=1, Dy=1, ticksize=-2pt 0, subticks=1]{->}(0,0)(-3.99,\ymin)(\xmax,\ymax)[{\footnotesize $x$},135][{\footnotesize $y$},-45]
\psplot[linewidth=1.2pt, plotpoints=200]{\xmin}{\xmax}{3*EXP(-x*x/2)}
\end{pspicture*}
\end{center}
\end{multicols}

\exo

La fonction $p$ est définie par la phrase suivante : \og à tout entier compris entre $100$ et $199$, la fonction $p$ associe son chiffre des dizaines \fg{}.
\begin{enumerate}
\item Si possible, déterminer les images par $p$ des nombres :
\[124\qpvq 256 \qpvq 162 \qpvq 103,7.\]
\item Si possible, déterminer un entier $n$ tel que $p(n)=5$.
\item Déterminer les antécédents de $3$ par $p$.
\end{enumerate} 

\exo

On considère la fonction $f$ définie sur $\R$ par $f(x)=x^3+5$, et on note $\cC_f$ sa courbe représentative dans un repère.
\begin{enumerate}
\item Calculer l'image de $10$ par $f$.

Que peut-on en déduire pour le point $A$ de coordonnées $\couple{10}{1005}$ ?
\item Calculer l'ordonnée du point $B$, appartenant à $\cC_f$ et ayant pour abscisse $-2$.
\end{enumerate}

\exo

On considère la fonction $g$ définie sur $\R$ par $g(x)=2x^3-3x+1$, et on note $\cC_g$ sa courbe représentative dans un repère.
\begin{enumerate}
\item Déterminer l'image de $2$ par $g$.
\item En déduire les coordonnées d'un point appartenant à la courbe $\cC_g$.
\item Déterminer les coordonnées d'un autre point appartenant à $\cC_g$.
\end{enumerate}

\exo

On considère la fonction $h$ définie sur $\R$ par $h(x)=5x+2$, et on note $\cC_h$ sa courbe représentative dans un repère.
\begin{enumerate}
\item Le point $E\couple{\frac{3}{5}}{4}$ appartient-il à $\cC_h$ ?
\item Calculer l'abscisse du point $F$, appartenant à $\cC_h$ et ayant pour ordonnée $0$.
\end{enumerate}

\exo

On considère la fonction $f$ définie sur $\R$ par $f(x)=3x^2+6$, et on note $\cC_f$ sa courbe représentative dans un repère.
\begin{enumerate}
\item Le point $A\couple{-1}{9}$ appartient-il à $\cC_f$ ?
\item Calculer l'ordonnée du point $B$ d'abscisse $4$ qui appartient à $\cC_f$.
\item Existe-t-il des points de $\cC_f$ dont l'ordonnée est égale à $33$ ?
\end{enumerate}

\exo

La fonction $f$ est définie sur $\intff{0}{5}$ par $f(x)=4-(x-3)^2$.
\begin{enumerate}
\item Construire un tableau de valeurs de la fonction $f$ avec un pas égal à $1$.
\item Dans un repère d'unités $1$ centimètre, placer six points appartenant à $\cC_f$.
\item Proposer un tracé de la courbe $\cC_f$.
\end{enumerate}


\exo

\vspace{-12pt}
\setlength{\columnsep}{1cm}
\setlength{\columnseprule}{0pt}
\begin{multicols}{2}
La courbe ci-contre est la représentation graphique d'une fonction $f$ définie sur $\intff{-4}{4}$.

Par lectures graphiques, déterminer l'ensemble de solutions de chaque équation :
\begin{enumerate}
\item $f(x)=1$ ;
\item $f(x)=2$ ;
\item $f(x)=0,5$ ;
\item $f(x)=3$.
\end{enumerate}

\begin{center}
\def\xmin{-5}
\def\xmax{5}
\def\ymin{-3}
\def\ymax{3}
\def\xunite{0.8cm}
\def\yunite{0.8cm}
\psset{unit=\xunite, xunit=\xunite, yunit=\yunite, algebraic=true, dimen=middle, dotstyle=*, dotsize=3pt 0, linewidth=0.8pt, arrowsize=3pt 2, arrowinset=0.25, comma=true}
\begin{pspicture*}(\xmin,\ymin)(\xmax,\ymax)
\psgrid[xunit=\xunite, yunit=\yunite, subgriddiv=0, gridlabels=0, gridcolor=lightgray](0,0)(\xmin,\ymin)(\xmax,\ymax)
\psaxes[labelFontSize=\scriptstyle, xAxis=true, yAxis=true, Dx=1, Dy=1, ticksize=-2pt 0, subticks=1]{->}(0,0)(-4.99,-2.99)(\xmax,\ymax)[{\footnotesize $x$},135][{\footnotesize $y$},-45]
\pscurve[linewidth=1.2pt](-4,2)(-1,0)(0.5,1)(3,-2)(4,-1)
\begin{footnotesize}
\uput{3pt}[000](4,-1){$\cC_f$}
\end{footnotesize}
\end{pspicture*}
\end{center}
\end{multicols}
\vspace{-6pt}

\exo

\vspace{-12pt}
\setlength{\columnsep}{1cm}
\setlength{\columnseprule}{0pt}
\begin{multicols}{2}
Les courbes ci-contre sont les représentations graphiques de trois  fonctions $f$, $g$ et $h$ définies sur $\intff{-4}{4}$.

Par lectures graphiques, déterminer l'ensemble de solutions de chaque équation :
\begin{enumerate}
\item $f(x)=g(x)$ ;
\item $f(x)=h(x)$ ;
\item $g(x)=h(x)$.
\end{enumerate}

\begin{center}
\def\xmin{-5}
\def\xmax{5}
\def\ymin{-4}
\def\ymax{3}
\def\xunite{0.8cm}
\def\yunite{0.8cm}
\psset{unit=\xunite, xunit=\xunite, yunit=\yunite, algebraic=true, dimen=middle, dotstyle=*, dotsize=3pt 0, linewidth=0.8pt, arrowsize=3pt 2, arrowinset=0.25, comma=true}
\begin{pspicture*}(\xmin,\ymin)(\xmax,\ymax)
\psgrid[xunit=\xunite, yunit=\yunite, subgriddiv=0, gridlabels=0, gridcolor=lightgray](0,0)(\xmin,\ymin)(\xmax,\ymax)
\psaxes[labelFontSize=\scriptstyle, xAxis=true, yAxis=true, Dx=1, Dy=1, ticksize=-2pt 0, subticks=1]{->}(0,0)(-4.99,-2.99)(\xmax,\ymax)[{\footnotesize $x$},135][{\footnotesize $y$},-45]
\psplot[linewidth=1.2pt]{-4}{4}{x^3/9-x^2/3-x+1}
\psplot[linewidth=1.2pt, linestyle=dashed]{-4}{4}{2-EXP(-0.462*x)}
\psplot[linewidth=1.2pt, linestyle=dotted]{-4}{4}{x^2/9-3}
\begin{footnotesize}
\uput{3pt}[000](4,1.8){$\cC_f$}
\uput{3pt}[135](-1.5,1.4){$\cC_g$}
\uput{3pt}[180](-4,-1.3){$\cC_h$}
\end{footnotesize}
\end{pspicture*}
\end{center}
\end{multicols}
\vspace{-6pt}

\exo

En reprenant le graphique de l'exercice \textbf{14}, déterminer l'ensemble de solutions de chaque inéquation :
\vspace{-12pt}
\setlength{\columnsep}{1cm}
\setlength{\columnseprule}{0pt}
\begin{multicols}{5}
\begin{enumerate}
\item $f(x)\geqslant 0$ ;
\item $f(x)<0$ ;
\item $f(x)<1$ ;
\item $f(x)\geqslant 1$ ;
\item $f(x)\leqslant 0,5$.
\end{enumerate}
\end{multicols}
\vspace{-6pt}

\exo

En reprenant le graphique de l'exercice \textbf{15}, déterminer l'ensemble de solutions de chaque inéquation :
\vspace{-12pt}
\setlength{\columnsep}{1cm}
\setlength{\columnseprule}{0pt}
\begin{multicols}{5}
\begin{enumerate}
\item $f(x)\leqslant g(x)$ ;
\item $f(x)>h(x)$ ;
\item $g(x)>h(x)$ ;
\item $g(x)\geqslant h(x)$.
\end{enumerate}
\end{multicols}
\vspace{-6pt}

\exo

Une fonction $f$ a les propriétés suivantes :
\begin{itemize}
\item $f$ est définie sur $\intff{0}{8}$ ;
\item l'équation $f(x)=3$ a pour solutions $1$ et $3$ ;
\item l'image de $0$ est $1$ ;
\item l'inéquation $f(x)\leqslant 0$ a pour ensemble de solutions $\intff{5}{7}$.
\end{itemize}
Tracer dans un repère une courbe qui pourrait représenter la fonction $f$.

\exo

La fonction $f$ est définie sur $\intff{-2}{2}$ par $f(x)=(3x+1)(5-x)$.

Sur la calculatrice, construire le tableau de valeurs de $f$ avec un pas égal à $0,5$.

\exo

Les membres d'un club d'astromodélisme ont construit une petite fusée : son altitude (en mètres) en fonction du temps $t$ (en secondes) est donnée par $h(t)=25t-5t^2$.

\begin{enumerate}
\item Quelle est l'altitude la fusée après $3$ secondes ?
\item Quelle est la durée du vol de la fusée ?
\item En utilisant la calculatrice, déterminer approximativement la durée pendant laquelle la fusée reste à plus de $10$ mètres d'altitude.
\end{enumerate}

\exo

\begin{minipage}[t]{7cm}
La fonction $f$ est définie sur $\intff{-2}{6}$ par :
\[f(x)=0,5x^2-2x-1.\]
Sa courbe représentative est tracée ci-contre.
\begin{enumerate}
\item Par lecture graphique, déterminer des valeurs approchées des solutions de l'équation $f(x)=1$.
\end{enumerate}
\end{minipage}
\hfill
\begin{minipage}[t]{10cm}
~

\vspace{-18pt}
\def\xmin{-3}
\def\xmax{7}
\def\ymin{-4}
\def\ymax{6}
\def\xunite{1.0cm}
\def\yunite{0.5cm}
\psset{unit=\xunite, xunit=\xunite, yunit=\yunite, algebraic=true, dimen=middle, dotstyle=*, dotsize=3pt 0, linewidth=0.8pt, arrowsize=3pt 2, arrowinset=0.25, comma=true}
\begin{pspicture*}(\xmin,\ymin)(\xmax,\ymax)
\psgrid[xunit=0.5cm, yunit=\yunite, subgriddiv=0, gridlabels=0, gridcolor=lightgray](0,0)(-6,\ymin)(14,\ymax)
\psaxes[showorigin=false,labelFontSize=\scriptstyle, xAxis=true, yAxis=true, Dx=2, Dy=2, ticksize=-2pt 0, subticks=1]{->}(0,0)(-2.99,-3.99)(\xmax,\ymax)[{\footnotesize $x$},135][{\footnotesize $y$},-45]
\psplot[linewidth=1.2pt, plotpoints=200]{-2}{6}{0.5*x^2-2*x-1}
\begin{footnotesize}
\uput{3pt}[135](5.5,4){$\cC_f$}
\end{footnotesize}
\end{pspicture*}
\end{minipage}

\medskip
\begin{enumerate}
\setcounter{enumi}{1}
\item On a construit un tableau de valeurs de la fonction $f$ :
\begin{center}
\setlength{\tabcolsep}{0.1cm}
\renewcommand{\arraystretch}{1.2}
\begin{tabular}{
|>{\centering\arraybackslash}m{2cm}
*{6}{|>{\centering\arraybackslash}m{1.5cm}}|}
\hline 
$x$ & $4,5$ & $4,6$ & $4,7$ & $4,8$ & $4,9$ & $5,0$ \\ 
\hline
$f(x)$ & $0,125$ & $0,380$ & $0,645$ & $0,920$ & $1,205$ & $1,500$ \\ 
\hline 
\end{tabular} 
\end{center}
\begin{enumerate}
\item En utilisant ces valeurs, donner un encadrement d'une des solutions de l'équation $f(x)=1$.
\item Quelle est l'amplitude de cet encadrement ?
\end{enumerate}
\item En utilisant la calculatrice, déterminer un encadrement d'amplitude $0,01$ de l'autre solution de cette équation.
\end{enumerate}

\exo

Une entreprise fabrique des pièces détachées pour automobiles. On note $x$ le nombre de pièces fabriquées au cours d'une journée. Le coût de production de $x$ pièces est noté $C(x)$ (en centaines d'euros).

D'autre part, chaque pièce est vendue $20$ euros. On note $R(x)$ la recette de l'entreprise pour $x$ pièces (en centaines d'euros).

Sur le graphique ci-dessous, on a tracé les courbes représentatives des fonctions $C$ et $R$ :

\begin{center}
\def\xmin{36}
\def\xmax{84}
\def\ymin{-2}
\def\ymax{20}
\def\xunite{0.25cm}
\def\yunite{0.25cm}
\psset{unit=\xunite, xunit=\xunite, yunit=\yunite, algebraic=true, dimen=middle, dotstyle=*, dotsize=3pt 0, linewidth=0.8pt, arrowsize=3pt 2, arrowinset=0.25, comma=true}
\begin{pspicture*}(\xmin,\ymin)(\xmax,\ymax)
\psgrid[xunit=\xunite, yunit=\yunite, subgriddiv=0, gridlabels=0, gridcolor=lightgray](0,0)(\xmin,\ymin)(\xmax,\ymax)
\psaxes[labelFontSize=\scriptstyle, xAxis=true, yAxis=true, Dx=5, Ox=40, Dy=2, ticksize=-2pt 0, subticks=1]{->}(40,0)(40,0)(\xmax,\ymax)[{\footnotesize $x$},135][{\footnotesize $y$},-45]
\pspolygon[fillcolor=black,fillstyle=solid,opacity=0.1](1.73,3)(1.73,5)(-1.73,5)(-1.73,3)
\psplot[linewidth=1.2pt, plotpoints=200]{40}{80}{0.01*x^2-0.79*x+17.48}
\psplot[linewidth=1.2pt, plotpoints=200]{40}{80}{0.2*x}
\psline[linewidth=1.2pt](-3,5)(3,5)
\psline[arrows=->,linewidth=1.2pt](0,0)(1,0)
\psline[arrows=->,linewidth=1.2pt](0,0)(0,1)
\begin{footnotesize}
\uput{3pt}[000](80,18.28){$\cC_C$}
\uput{3pt}[000](80,16){$\cC_R$}
\uput{3pt}[90](2.7,5){$\Delta$}
\uput{3pt}[270](1,0){$1$}
\uput{3pt}[180](0,1){$1$}
\end{footnotesize}
\end{pspicture*}
\end{center}

\begin{enumerate}
\item À l'aide du graphique :
\begin{enumerate}
\item déterminer le coût de production de $50$ pièces ;
\item déterminer le nombre de pièces que l'entreprise peut produire pour un coût égal à \np{1000} euros.
\end{enumerate}
\item
\begin{enumerate}
\item En utilisant les données de l'énoncé, déterminer l'expression de la recette $R(x)$ pour $x$ pièces.
\item Vérifier que cette expression est cohérente avec le tracé de $\cC_R$.
\end{enumerate}
\item Pour $x$ pièces, le bénéfice $B(x)$ de l'entreprise est donné par $B(x)=R(x)-C(x)$. S'il est négatif, ce \og bénéfice \fg{} est donc une perte.

Combien de pièces l'entreprise doit-elle fabriquer et vendre pour réaliser un bénéfice positif ?
\end{enumerate}

